\begin{helper}
Let $D>0$ be an integer, let $\mathcal H\in H(3,2,D)$ have finitely
many edge indices (with no finiteness assumption on its vertex set), and
let $e$ be an edge index. There exist natural numbers $m,\ell$ with
$m\le 1+3(D-1)$ and $\ell\le 3m$, a hypergraph
$\mathcal F\in H(3,2,D)$ on vertex set $\operatorname{Fin}(\ell)$ and
edge-index set $\operatorname{Fin}(m)$, and an index $f$, such that
\[
b_{L(\mathcal H),3D}(e)=b_{L(\mathcal F),3D}(f).
\]
\end{helper}
\begin{proof}
By \noderef{HypergraphClass}, $\mathcal H$ also belongs to $H(3,3,D)$:
its intersection bound two implies the bound three. Apply
\noderef{KUniformKSimpleLocalTruncationPreservesBParameter} with $k=3$.
Write $\mathcal K,a$ for the resulting truncation. Its edge-index
bijection identifies its indices with $e$ and the line-graph neighbors of
$e$; its vertex bijection identifies its vertices with those lying in
these retained edges. Its incidence equivalence identifies each local
edge with its original edge, and its local $b$-value equals the original
value. For distinct local indices, the vertex map injects their
intersection into the intersection of the corresponding distinct original
edges. Thus their intersection has size at most two. The truncation's
uniformity and degree bounds, together with this observation, show that
$\mathcal K\in H(3,2,D)$.

Every index other than $a$ intersects $a$, by the incidence equivalence
and the edge-index bijection. For each $x\in\mathcal K(a)$ let $A_x$ be
the set of indices other than $a$ containing $x$. Its size is
$\deg_{\mathcal K}(x)-1\le D-1$. The union of these three sets contains
every index other than $a$, so, writing $m$ for the number of local
indices, $m-1\le 3(D-1)$. Let $S$ be the union of all local edges and
put $\ell=|S|$. The union bound and uniformity give $\ell\le 3m$.

Choose bijections from $\operatorname{Fin}(\ell)$ to $S$ and from the
local edge indices to $\operatorname{Fin}(m)$. Define $\mathcal F$ by
transporting incidence under these bijections, and take $f$ to be the
image of $a$. Each edge and each pairwise intersection is in bijection
with its counterpart in $\mathcal K$. Each vertex's incident-index set
is likewise in bijection with the old incident-index set. Hence edge
sizes, intersection sizes, and vertex degrees are preserved, giving
$\mathcal F\in H(3,2,D)$. By \noderef{LineGraphOfHypergraph}, the
edge-index bijection preserves and reflects line-graph adjacency. It
therefore maps the neighborhood of $a$ bijectively onto that of $f$ and
preserves adjacency within it. Apply
\noderef{LocalBParameterNeighborhoodEmbeddingTransport} to obtain
equality of their local $b$-values. Combine this equality with the
truncation equality to finish. No additional edge indices are introduced.
\end{proof}
